\chapter{The Unstructured Mesh}\label{ch:mesh}

This chapter describes the static geometry on which the engine operates: the
unstructured mesh, the quantities a model stores on it, and the conventions under which each source ---
FESOM2, FESOM~1.4, FESOM-C, HEALPix products and structured NEMO grids --- is read into the \fn{Mesh} structure of
\cref{sec:arch-data}. We use the triangular FESOM2 \emph{core2} mesh as the worked example and describe
the other sources where they differ. The localization search of \cref{ch:localization} and every
interpolation stencil of \crefrange{ch:hvel}{ch:sampling} are expressed in terms of
the entities introduced here, so we also fix the vocabulary --- node, element,
interface, layer --- used throughout the manual.

\section{Nodes, elements, and interfaces}\label{sec:mesh-anatomy}

The horizontal mesh is a tessellation of the ocean surface: a set of $N$
\emph{nodes} joined into $E$ \emph{elements}. \index{node}\index{element} In FESOM2
the elements are triangles; in FESOM-C they are mixed, each element being a triangle
or a quadrilateral; structured grids and the HEALPix pixel mesh produce quadrilaterals only, and the
HEALPix dual mesh quadrilaterals with a few triangles (\cref{sec:mesh-sources}). A triangle in a mesh with room for four vertices is stored with its fourth slot
set to $-1$.
For the FESOM2 \emph{core2} configuration used as the running example the elements are
triangular, with $N=\num{126858}$ nodes and $E=\num{244659}$ elements. Each element
stores the indices of its vertices --- three for a triangle, four for a quadrilateral
--- and the engine precomputes, for each node, the list of elements that contain it
(its \emph{one-ring}), which the localization search relies upon
(\cref{ch:localization}).

The model does not store all fields at the same location, and the distinction is
fundamental to interpolation (\cref{fig:mesh-anatomy}). \emph{Node-layout} fields
are defined at the vertices: the FESOM2 horizontal velocity components \conf{unod}
and \conf{vnod}, the sea-surface height, and the sampled scalars (temperature,
salinity) are all node quantities. \emph{Element-layout} fields are defined at the
element centroids: the alternative monthly-mean FESOM2 velocities \conf{u} and
\conf{v} are element quantities. A node-layout field carries one value per vertex
and is interpolated by the barycentric scheme of \cref{ch:hvel}; an element-layout
field carries one value per element and is interpolated by either least-squares scheme or,
through area-weighted vertex values, the barycentric scheme of the same chapter. The \fn{VelocityField} records which layout applies in
its \code{uv\_layout} attribute, and the engine selects the matching stencil
accordingly.

\begin{figure}[t]
\centering
\begin{tikzpicture}[scale=1.25]
  \coordinate (C)  at (2.00,1.55);
  \coordinate (p1) at (0.35,1.50);
  \coordinate (p2) at (1.15,0.20);
  \coordinate (p3) at (2.95,0.35);
  \coordinate (p4) at (3.70,1.70);
  \coordinate (p5) at (2.75,2.95);
  \coordinate (p6) at (1.00,2.80);
  % the one-ring of triangles around node C
  \foreach \a/\b in {p1/p2,p2/p3,p3/p4,p4/p5,p5/p6,p6/p1}
     \fill[meshface] (C)--(\a)--(\b)--cycle;
  % highlight one element
  \fill[hilface] (C)--(p2)--(p3)--cycle;
  \foreach \a/\b in {p1/p2,p2/p3,p3/p4,p4/p5,p5/p6,p6/p1}{
     \draw[meshedge] (C)--(\a); \draw[meshedge] (\a)--(\b);}
  % element centroid of (C,p2,p3)
  \node[elemcentre] (E) at ($ 0.3333*(C) + 0.3333*(p2) + 0.3333*(p3) $) {};
  % nodes
  \foreach \p in {C,p1,p2,p3,p4,p5,p6} \node[meshnode] at (\p) {};
  % labels
  \node[nodelbl,above right=1pt of C] {node $n$: $(\conf{unod},\conf{vnod})$, $T$, $S$};
  \node[elemlbl,below right=0pt of E] {element centre: $(u_e,v_e)$};
  \node[font=\footnotesize,gray!60!black] at (2.05,0.62) {highlighted element};
\end{tikzpicture}
\caption{Anatomy of the horizontal mesh. Triangle vertices are \emph{nodes}
(filled circles), where node-layout velocities and the scalar fields live; triangle
centroids are \emph{elements} (the orange square), where element-layout velocities
live. The six shaded triangles sharing the central node $n$ form its one-ring, the
adjacency the localization search exploits (\cref{ch:localization}).}
\label{fig:mesh-anatomy}
\end{figure}

\section{Vertical structure: interfaces and layers}\label{sec:mesh-vertical}

In the vertical, the water column is divided by $K$ \emph{interfaces} into $K-1$
\emph{layers} (\cref{fig:mesh-column}). \index{interface}\index{layer} The
interfaces are the level boundaries; the layers are the volumes between them. The
distinction matters because, again, different fields live in different places: the
vertical velocity $\Wvel$ is defined \emph{at the interfaces}, whereas the
three-dimensional horizontal velocity and the scalar fields are defined \emph{at
the mid-layers}. FESOM~1.4 is the exception: it stores every field on the levels, so its horizontal
velocity shares the interfaces with $\Wvel$ (\cref{sec:mesh-fesom14}). For the \emph{core2} mesh there are $48$ interfaces and hence $47$
layers, spanning depths from $0$ to approximately \SI{6250}{\metre}. The
configuration key \conf{max\_levels} retains only the first $N$ interfaces from the
surface downward (at least two in three dimensions, on either vertical family), which both saves memory
and caps the depth a particle can reach; a run that follows water to the sea floor keeps the full
column.

\begin{figure}[t]
\centering
\begin{tikzpicture}[scale=1.0]
  \def\xa{0} \def\xb{3.4}
  % interface depths (increasing layer thickness with depth)
  \foreach \y/\lab in {0/{$z_0=0$ (surface)},-0.5/{$z_1$},-1.1/{$z_2$},-1.8/{$z_3$},-2.6/{$z_4$},-3.5/{$z_5$}}{
     \draw[interface] (\xa,\y)--(\xb,\y);
     \node[interfacept] at (1.1,\y) {};
     \node[font=\footnotesize,left] at (\xa,\y) {\lab};
  }
  \draw[interface] (\xa,-4.5)--(\xb,-4.5);
  \node[interfacept] at (1.1,-4.5) {};
  \node[font=\footnotesize,left] at (\xa,-4.5) {$z_6$};
  % mid-layer markers
  \foreach \y in {-0.25,-0.8,-1.45,-2.2,-3.05,-4.0}
     \node[font=\footnotesize,lagelem!70!black] at (2.55,\y) {$\times$};
  % local sea floor
  \draw[dashed,laggreen,line width=0.9pt] (\xa,-4.05)--(\xb,-4.05);
  \node[font=\footnotesize,laggreen,right] at (\xb,-4.05) {local floor (\conf{zbar\_n\_bottom})};
  \fill[pattern=north east lines,pattern color=gray!45] (\xa,-4.05) rectangle (\xb,-4.55);
  % depth axis
  \draw[-{Stealth[length=5pt]},gray] (-1.65,0.0)--(-1.65,-4.3);
  \node[gray,rotate=90,anchor=south,font=\footnotesize] at (-1.85,-2.1) {depth $\depx$ (positive down)};
  % field-location annotations
  \node[font=\footnotesize,interfacept,label={[font=\footnotesize]right:{interfaces: $\Wvel$}}] at (1.1,-2.05) {};
  \node[font=\footnotesize,lagelem!70!black] at (2.55,-2.62) {};
  \node[font=\footnotesize,right] at (2.75,-1.45) {mid-layers: $u,v$ (3-D), $T$, $S$};
\end{tikzpicture}
\caption{Anatomy of the vertical column. The $K$ interfaces (grey lines, open
circles) bound $K-1$ mid-layers (orange crosses). The vertical velocity $\Wvel$ is
stored at the interfaces; the three-dimensional horizontal velocity and the scalar
fields are stored at the mid-layers. The local sea floor is given per node by
\conf{zbar\_n\_bottom}; layers below it do not exist for that node and a particle
cannot descend past it (\cref{ch:vertical}).}
\label{fig:mesh-column}
\end{figure}

The engine treats FESOM2 levels as fixed depths (a \emph{z-level} vertical coordinate), the same
everywhere. In FESOM2 itself the interfaces follow the free surface (the zstar coordinate), and the
recorded depth is therefore the reference depth, within $|\eta|$ of the physical one
(\cref{sec:vert-wsign}).
The alternative supported family, used by the FESOM-C model, is a
terrain-following \emph{sigma} coordinate in which the interface depths are a
fixed fraction of the local water-column thickness and therefore vary in space and
with the sea-surface height. The two families differ only in how an interface depth
is computed; we defer the sigma geometry to \cref{ch:vertical}. FESOM~1.4, HEALPix products and NEMO
grids use z-levels.

\section{Reading a mesh file}\label{sec:mesh-conventions}

The FESOM2 reader accepts three mesh sources: a mesh-diagnostics NetCDF file, an ASCII mesh directory, and
a Zarr store written by the region-extraction tool. It resolves the quantities of a diagnostics file by
variable \emph{name} rather than by dimension, which lets it read three naming conventions --- node
coordinates as \code{lon}/\code{lat} in degrees or packed as one \code{nodes} array in radians, and the
connectivity as \code{elements}, \code{elem} or \code{face\_nodes}, oriented to $(E,3)$ by its axis of
length three --- and applies a small number of unit and indexing conversions (\cref{tab:mesh-vars}). Packed
coordinates are converted to degrees, trusting an explicit \code{units} attribute when present and otherwise
inferring radians when the longitude magnitude does not exceed $2\pi$ and the latitude magnitude does not
exceed $\pi/2$. The connectivity is shifted to zero-based indexing by its \code{start\_index} attribute
(one, as FESOM2 writes it from Fortran, when absent). Depth-like quantities --- the interface levels and the
per-node sea-floor depth --- are stored with a negative-downward sign and are converted to the positive-down
convention by taking absolute values. Element centres are computed from the vertices, with the longitude
averaged on the circle so that an element across the dateline gets a centre inside it. The model's element
areas are kept when the file provides them.

\begin{table}[t]
\centering
\caption{Quantities read from the FESOM2 mesh-diagnostic file and their conversion
to the \fn{Mesh} structure.}
\label{tab:mesh-vars}
\small
\begin{tabularx}{\linewidth}{@{}l l X@{}}
\toprule
File variable & \fn{Mesh} member & Conversion / convention \\
\midrule
\code{lon}, \code{lat} or \code{nodes} $(2,N)$ & \code{nodes\_lon/lat} & degrees; packed radians $\rightarrow$ degrees \\
\code{elements}, \code{elem} or \code{face\_nodes} & \code{elem\_nodes}    & oriented to $(E,3)$; one-based $\rightarrow$ zero-based \\
\code{nz} or \code{zbar}    & \code{sigma}, \code{z\_levels} & $|\cdot|$, positive-down interface depths \\
\code{zbar\_n\_bottom}      & \code{depth}          & $|\cdot|$, per-node sea floor (positive-down) \\
\code{elem\_area}           & element areas         & used directly when present \\
\code{nod\_in\_elem2D} (+ count) & \code{node\_to\_elems} & one-based $\rightarrow$ zero-based; rebuilt if degenerate (\cref{sec:mesh-nodeelem}); required in a diagnostics file \\
\code{zbar\_e\_bottom}, \code{nlevels} & element bottoms (metadata) & each element's bed and last whole layer, when present (\cref{sec:mesh-elem-bottoms}) \\
\code{ulevels}             & element first levels (metadata) & above 1 under an ice-shelf cavity \\
\bottomrule
\end{tabularx}
\end{table}

\codeanchor{Mesh-variable resolution and conversions are in \fn{\_resolve\_mesh\_vars},
\fn{\_load\_mesh\_nc} and \fn{\_compute\_elem\_centres} (\srcfile{readers/fesom2.py}); the \fn{Mesh}
dataclass and its \code{z\_levels} property are in \srcfile{core/mesh.py}.}

\section{The node-to-element table and its reconstruction}\label{sec:mesh-nodeelem}

The localization search requires, for each node, the list of elements that contain
it; this is the \code{nod\_in\_elem2D} table in the mesh file. On the \emph{core2}
mesh this table is effectively empty: of the $\num{733977}$ node--element incidences it should hold
(three per element) only $54$ are present. \index{nod\_in\_elem2D@\code{nod\_in\_elem2D}} The consequence, if the
table were used as read, is severe and silent. After the routine one-based to
zero-based shift, every unpopulated entry becomes $-1$, so that $\num{126852}$ of
the $\num{126858}$ nodes report no adjacent elements. The containing-element search
of \cref{ch:localization} would then fail for every particle on the first step,
flag each as having left the domain, and terminate normally --- producing an output
file in which no particle ever moved, with no error raised.

To prevent this, the reader detects a degenerate table and rebuilds it from the
element connectivity, which is intact. The detector flags the table as degenerate on
any of three independent signals: more than half of the nodes report no first
element; an entry strictly inside a node's declared neighbour count is empty (the
table and the count array disagree); or a neighbour count is negative or exceeds the
maximum connectivity. On a positive verdict the engine prints a notice and
reconstructs the node-to-element adjacency by a two-pass count-and-fill over the
elements, after which the localization search behaves correctly. We note this
behaviour explicitly because it is a precondition for any meaningful run on the
\emph{core2} mesh, and because its failure mode --- a completed run with no motion
--- is easy to misread as a physical result rather than a missing table.

\codeanchor{The degeneracy detector \fn{\_nod\_in\_elem2d\_is\_degenerate} and the
rebuild call are in \srcfile{readers/fesom2.py}; the reconstruction
\fn{build\_node\_to\_elems} (a two-pass count-and-fill) is in \srcfile{core/mesh.py}.}

\section{Element bottoms}\label{sec:mesh-elem-bottoms}

FESOM2 fixes the number of levels of each element first and gives a node the deepest of the elements
around it, so the per-node floor \code{zbar\_n\_bottom} tells where the deepest element at a node ends,
not where a given element ends. The depth-linear reading of element-layout velocities
(\cref{sec:vert-uvdepth}) needs the bed of every element, and the FESOM2 reader takes it from the mesh
diagnostic --- NetCDF or Zarr --- when the diagnostic holds the element level variables: \code{nlevels}
(one-based; the element has $\code{nlevels}-1$ wet layers), \code{zbar\_e\_bottom} (its bed, negative
down) and \code{ulevels} (its first level, above 1 under an ice-shelf cavity). From them it declares, in
file element order, the bed $d_e$ and the last whole layer $D_e$ of each element,
\begin{equation}\label{eq:mesh-elem-bed}
  d_e = \begin{cases} |\code{zbar\_e\_bottom}_e| & \text{when present},\\
                      z_{\code{nlevels}_e-1} & \text{otherwise},\end{cases}
  \qquad
  D_e = \max\{k : z_{k+1}\le d_e\}\quad(-1 \text{ when } z_1>d_e),
\end{equation}
that is, the last layer whose lower level lies at or above the bed. A bottom layer that the bed cuts ---
a FESOM2 \emph{partial bottom cell} --- is therefore not counted, and the depth-linear rule never reads it,
the same rule as the node table of \cref{eq:vert-uv-last}. The element first levels are kept as well.
A variable read here that is not one finite value per element, or that holds levels outside $1,\dots,K$,
refuses the mesh at load --- in layer runs too --- with a message that names the file and the fix, and
the loader prints one line saying whether the bottoms came from the file. A mesh without either
variable --- an ASCII mesh directory, a diagnostic that lacks them, a Zarr store without a level table, or
a region extracted before the extraction tool kept them --- declares no bottoms, and
\conf{mode.uv\_vertical: linear} on its element velocities is refused before the output file exists;
\code{ulevels} alone declares nothing. A mesh with \code{ulevels}${}>1$ anywhere is refused for the
same option: ice-shelf cavities are not supported in the code yet and are planned. The region-extraction
tool \code{utils/extract\_fesom\_region.py} writes \code{nlevels}, \code{ulevels} and
\code{zbar\_e\_bottom} of the kept elements, and \code{nlevels\_nod2D} of the kept nodes, whenever the
source has them, with the values and signs of the source; a subset written before it did so must be
extracted again to use the option on element data.

\codeanchor{\raggedright \fn{\_elem\_level\_metadata} in \srcfile{readers/fesom2.py}, called by the NetCDF and the Zarr
mesh loaders, declares \code{mesh.metadata["elem\_bed"]}, \code{["uv\_data\_last\_elem"]},
\code{["elem\_ulevels"]} and \code{["elem\_levels\_source"]}; the element table that reads them is
\fn{depth\_points} with \code{layout="element"} (\srcfile{numerics/vertical\_uv.py}); the subset writer is
\srcfile{utils/extract\_fesom\_region.py}.}


\section{Other mesh sources}\label{sec:mesh-sources}

The other readers produce the same \fn{Mesh} from different inputs (\cref{tab:mesh-sources,fig:mesh-sources});
the FESOM~1.4 reader has a section of its own (\cref{sec:mesh-fesom14}).

\begin{table}[t]
\centering\footnotesize\setlength{\tabcolsep}{4pt}
\begin{tabular}{@{}>{\raggedright\arraybackslash}p{23mm}>{\raggedright\arraybackslash}p{29mm}>{\raggedright\arraybackslash}p{29mm}>{\raggedright\arraybackslash}p{17mm}>{\raggedright\arraybackslash}p{16mm}>{\raggedright\arraybackslash}p{14mm}@{}}
\toprule
source & nodes & elements & velocity layout & vertical & radius (m) \\
\midrule
FESOM2 & model nodes & triangles & node or element & z-level & \num{6367500} \\
FESOM~1.4 & model nodes & triangles, from the mesh files or the cell polygons & node, on the levels & z-level & \num{6367500} \\
FESOM-C & model nodes & triangles, quadrilaterals & element (node when stored so) & sigma & \num{6400000} \\
HEALPix, dual & wet pixel centres & corner cells: quadrilaterals, some triangles & node & z-level & \num{6371000} \\
HEALPix, pixel & pixel corners & wet pixels (quadrilaterals) & element & two-dim. & \num{6371000} \\
NEMO & wet T-points & quadrilaterals of four wet T-points & node & two-dim. & \num{6371000} \\
\bottomrule
\end{tabular}
\caption{Mesh sources. The radius is the default of \conf{mesh.r\_earth}; it enters the fallback
computation of element areas and, without effect on the values, the barycentric projection.}
\label{tab:mesh-sources}
\end{table}

\paragraph{FESOM-C.} The FESOM-C reader takes the node coordinates, depths and vertex lists from the
model's single output file, together with the element centres when present and the sigma levels. Its
connectivity is sanitised on load: indices are shifted to zero-based, empty and repeated vertices are
removed, and each row is compacted, so a triangle stored as a quadrilateral with a repeated vertex
becomes $[a,b,c,-1]$.

\paragraph{HEALPix, dual mesh.} A HEALPix product stores its fields at the centres of equal-area pixels.
The default \emph{dual} mesh takes the wet pixel centres as nodes and builds one element around each pixel
corner, joining the pixels that meet there: four pixels give a quadrilateral, and at the eight points of
the sphere where only three pixels meet, a triangle. The cells are built from the pixel neighbour
relations in integer arithmetic, each kept once; an all-wet sphere must give exactly two more elements than
pixels, which the builder checks. A pixel is wet when the velocity at the first record is finite on the
tracked level, restricted by \conf{mesh.crop\_bbox} and by a bathymetry variable when the store has one.
The velocities live at the nodes, so the dual mesh needs \conf{uv\_interp: barycentric}.

\paragraph{HEALPix, pixel mesh.} With \conf{mesh.healpix\_geometry: pixel} the elements are the wet pixels
themselves, quadrilaterals whose nodes are the pixel corners; the velocities live at the element centres,
so both least-squares schemes apply. A quadrilateral with straight edges approximates a pixel with curved
edges, so near an edge a point may be assigned to the adjacent pixel. The pixel mesh is two-dimensional;
the eight pixels that touch a pole are degenerate in the geographic frame and produce a warning.

\paragraph{Structured grids and NEMO.} A logically rectangular grid becomes a mesh of quadrilaterals whose
corners are four adjacent wet cell centres (T-points); a quadrilateral with a land corner is dropped, which
draws the coastline, and a wet T-point that belongs to no quadrilateral is left out, since it could only be
a dead end. The NEMO reader averages the C-grid velocities at the cell faces to the T-points, setting
missing face values to zero first, which gives node-layout velocities. Its first version is two-dimensional
and requires a grid whose axes are within $2^\circ$ of east and north.

\begin{figure}[t]
\centering
\begin{tikzpicture}[scale=0.9]
  % (a) HEALPix dual: pixel centres as nodes, a quad around a 4-valent corner, a triangle around a 3-valent one
  \begin{scope}
    \foreach \x/\y in {0/0,1.4/0,0/1.4,1.4/1.4} \draw[gray!50] (\x-0.7,\y-0.7) rectangle (\x+0.7,\y+0.7);
    \fill[hilface] (0,0)--(1.4,0)--(1.4,1.4)--(0,1.4)--cycle;
    \draw[meshedge] (0,0)--(1.4,0)--(1.4,1.4)--(0,1.4)--cycle;
    \foreach \x/\y in {0/0,1.4/0,0/1.4,1.4/1.4} \node[meshnode] at (\x,\y) {};
    \node[font=\footnotesize,align=center,text width=40mm] at (0.7,-1.35) {(a) dual: nodes at pixel\\centres, cell at a corner};
  \end{scope}
  % (b) pixel mesh: pixel as element, corners as nodes
  \begin{scope}[xshift=5.2cm]
    \fill[hilface] (-0.7,-0.7) rectangle (0.7,0.7);
    \draw[meshedge] (-0.7,-0.7) rectangle (0.7,0.7);
    \foreach \x/\y in {-0.7/-0.7,0.7/-0.7,0.7/0.7,-0.7/0.7} \node[meshnode] at (\x,\y) {};
    \node[elemcentre] at (0,0) {};
    \node[font=\footnotesize,align=center,text width=40mm] at (0,-1.35) {(b) pixel: element = pixel,\\velocity at its centre};
  \end{scope}
  % (c) structured grid: T-point quads, one land cell
  \begin{scope}[xshift=8.6cm,yshift=-0.7cm]
    \foreach \i in {0,1,2} \foreach \j in {0,1,2} {
      \ifnum\i=2 \ifnum\j=2 \else \node[meshnode] at (\i*0.8,\j*0.8) {};\fi \else \node[meshnode] at (\i*0.8,\j*0.8) {};\fi }
    \fill[hilface] (0,0) rectangle (0.8,0.8);
    \fill[hilface] (0.8,0) rectangle (1.6,0.8);
    \fill[hilface] (0,0.8) rectangle (0.8,1.6);
    \draw[meshedge] (0,0) rectangle (0.8,0.8) (0.8,0) rectangle (1.6,0.8) (0,0.8) rectangle (0.8,1.6);
    \node[rectangle,fill=gray!50,inner sep=2.2pt] at (1.6,1.6) {};
    \node[font=\footnotesize,anchor=west] at (1.75,1.6) {land};
    \node[font=\footnotesize,align=center,text width=40mm] at (0.8,-0.65) {(c) structured: quads of\\four wet T-points};
  \end{scope}
\end{tikzpicture}
\caption{Three mesh sources on one \fn{Mesh}. (a) The HEALPix dual mesh joins the four pixel centres around
a pixel corner into a quadrilateral (three centres give a triangle). (b) The pixel mesh makes each wet pixel
an element, with nodes at its corners and the velocity at its centre. (c) A structured grid keeps the
quadrilaterals whose four T-points are wet; one land point removes the quadrilateral in the corner.}
\label{fig:mesh-sources}
\end{figure}

\codeanchor{\fn{\_sanitize\_connectivity} and \fn{FesomCReader.load\_mesh} (\srcfile{readers/fesomc.py});
\fn{build\_healpix\_dual\_mesh}, \fn{build\_healpix\_pixel\_mesh} and \fn{HealpixReader.load\_mesh}
(\srcfile{readers/healpix.py}); \fn{build\_structured\_mesh} and \fn{c\_grid\_to\_t}
(\srcfile{readers/structured.py}); the NEMO adapter in \srcfile{readers/nemo.py}.}

\section{The FESOM~1.4 mesh}\label{sec:mesh-fesom14}

FESOM~1.4 \citep{wang2014fesom14} is a finite-element model on a triangular surface mesh that is the same
at every level. Every surface node has a column of three-dimensional nodes, one per level, down to its
bed, and all prognostic fields --- $u$, $v$, $\Wvel$, temperature and salinity --- are continuous,
piecewise-linear (P1) functions on the tetrahedra between those nodes: they live \emph{on the levels at the
nodes}. The reader therefore returns node-layout fields on the levels (\cref{sec:vert-uvdepth}), and it
builds the \fn{Mesh} from either of two sources, decided by what \conf{mesh.meshpath} is.

\paragraph{ASCII mesh directory.} A directory that holds \code{nod3d.out} is read as the model's mesh
files: \code{nod2d.out} (the node coordinates and flags), \code{elem2d.out} (the triangles, one-based),
\code{nod3d.out} (every three-dimensional node with its depth, numbered level by level) and
\code{aux3d.out} in FESOM~1.4's own format --- the number of levels, the table of the three-dimensional
node at each level of each column (a value $\le0$ marks a dry slot) and the surface node of every
three-dimensional node; its last block, on the three-dimensional elements, is not parsed. This
\code{aux3d.out} is not FESOM2's, and the FESOM2 reader would misread it. Each file is read as one stream
of numbers, as a Fortran list-directed read sees it, so the spread of the records over lines does not
matter; Fortran \code{D} exponents are accepted, and an integer field takes any number whose value is an
integer. The levels are the depths of \code{nod3d.out} in file order --- every change of value starts a new
level --- stored positive down, with the surface exactly at zero. Before anything is returned the files
are checked against each other: indices in range, every three-dimensional node in exactly one slot, the
wet slots a prefix of each column, and every slot's surface node and depth consistent with the node
tables. A failure names the file and the first offending index.

A FESOM~1.4 run may store its mesh in the rotated frame of the model. The key \conf{mesh.input\_rotation}
takes the Euler angles $[\alpha,\beta,\gamma]$ of the run's namelist, in degrees, and rotates the node
coordinates back to geographic on load, with the convention of \cref{sec:frame-rotation}. The angles are
never guessed, since a rotated mesh looks like a geographic one; the reader warns when a global ASCII mesh
(more than $350^\circ$ of longitude) without a node north of $85^\circ$\,N is loaded without them, because a
geographic FESOM mesh reaches the Arctic. The key is refused for a mesh from cell polygons, whose
coordinates are geographic, and for every other mesh type; it is independent of \conf{mesh.rotation}, the
frame in which the engine tracks (\cref{ch:frame}).

\paragraph{Cell polygons of level-wise data.} Published FESOM~1.4 output, such as the CMIP6 files, holds
no element table. It holds, next to the node coordinates \code{lon}/\code{lat}, each node's
\emph{median-dual cell} as a polygon (\code{lon\_bnds}/\code{lat\_bnds}, padded to a fixed vertex count by
repeating the last vertex). An interior node's polygon alternates the midpoints of its edges and the
centroids of its triangles; a coastal node's polygon also holds the midpoints of its boundary edges and the
node itself. A polygon vertex shared by exactly three nodes is therefore the centroid of their triangle
(\cref{fig:mesh-fesom14}), and the reader recovers the triangles exactly from that rule. The polygon
vertices of all nodes are clustered on the unit sphere within $10^{-6}$ degrees (about \SI{0.1}{\metre});
distinct clusters must lie at least a thousand times farther apart, so that no point split by rounding can
pass for two. At a \emph{pinch node}, whose triangles form several fans that touch only at the node, the
published polygon lists one fan only; a centroid listed by just two nodes then names its third node, and
that triangle is recovered as well. The result is validated before it is used: a shared point within
\SI{1}{\centi\metre} of the spherical centroid of its three nodes, every edge in one or two triangles with a
polygon vertex at its midpoint, the Euler characteristic $V-E+F=2C-L$ of $C$ edge-connected components with
$L$ boundary loops, and an accounting check that every distinct polygon vertex is a node, the midpoint of a
recovered edge or the centroid of a recovered triangle. The levels are the file's depth coordinate (finite,
starting at the surface, increasing), and the wet mask comes from the first record of $u$: a value counts
as data when it is finite and smaller than $10^{29}$ in magnitude, and the data of every column must be a
prefix of its levels, starting at the surface. When a file holds \code{deptho} (for instance the published
fixed-field file in the same directory), its value must equal the depth of each column's last level with
data within \SI{0.01}{\metre}. This is the only guard against level-wise files that write zero rather than
a fill value below the bed, whose first record would look wet everywhere; without \code{deptho} that case
cannot be detected.

\begin{figure}[t]
\centering
\begin{tikzpicture}[scale=1.15]
  \coordinate (A) at (0,0); \coordinate (B) at (4,0); \coordinate (C) at (1.6,2.8);
  \coordinate (G) at (1.8667,0.9333);
  \coordinate (Mab) at (2,0); \coordinate (Mbc) at (2.8,1.4); \coordinate (Mca) at (0.8,1.4);
  \fill[lagnode!14] (A)--(Mab)--(G)--(Mca)--cycle;
  \fill[lagelem!14] (B)--(Mbc)--(G)--(Mab)--cycle;
  \fill[laggreen!14] (C)--(Mca)--(G)--(Mbc)--cycle;
  \draw[meshedge] (A)--(B)--(C)--cycle;
  \draw[lagnode!80!black,densely dashed,line width=0.7pt] (Mab)--(G) (Mbc)--(G) (Mca)--(G);
  \foreach \p in {A,B,C} \node[meshnode] at (\p) {};
  \foreach \p in {Mab,Mbc,Mca} \node[circle,fill=black,inner sep=0pt,minimum size=3pt] at (\p) {};
  \node[elemcentre] at (G) {};
  \node[nodelbl,below left] at (A) {$A$};
  \node[nodelbl,below right] at (B) {$B$};
  \node[nodelbl,above] at (C) {$C$};
  \node[font=\footnotesize,anchor=west,align=left] at (4.3,1.55) {$G$: a vertex of the polygons of $A$, $B$\\and $C$ --- the centroid of triangle $ABC$};
  \draw[gray!70,-{Stealth[length=4pt]}] (4.25,1.55)..controls (3.2,1.4)..(2.0,1.0);
  \node[font=\footnotesize,anchor=west,align=left] at (4.3,0.35) {edge midpoint: a vertex of two\\polygons --- the edge $AB$};
  \draw[gray!70,-{Stealth[length=4pt]}] (4.25,0.35)..controls (3.0,0.4)..(2.08,0.06);
\end{tikzpicture}
\caption{Triangle recovery from FESOM~1.4 cell polygons. The median-dual cell of each node (shaded, one
colour per node) is bounded by the dashed segments joining the edge midpoints (dots) to the triangle
centroid $G$ (square). Within one triangle each node's cell contributes the quadrilateral node, midpoint,
centroid, midpoint, so $G$ is the only point that the polygons of exactly three nodes share: every cluster
of polygon vertices owned by three nodes is a triangle of the mesh.}
\label{fig:mesh-fesom14}
\end{figure}

\paragraph{The mesh in either case.} \code{mesh.sigma} holds the full level table, positive down, and
\code{mesh.depth} the depth of each column's last wet level, exactly. The reader keeps the static wet mask
of every column (against which every record of every variable is later compared, \cref{sec:input-reader}),
the geographic node coordinates, and the index of each column's last wet level, which it declares as the
data prefix of the depth-linear rule ($D_i$ in \cref{eq:vert-uv-last}). Element centres are circular means
of the vertices, as for FESOM2, and the default radius is \SI{6367500}{\metre}, FESOM~1.4's value. A mesh
whose columns do not all start at the surface --- a mesh with ice-shelf cavities --- is refused by the
reader: an ASCII mesh whose top wet node lies below $z=0$, or level-wise data with a column that holds no
surface value. Ice-shelf cavities are not supported in the code yet and are planned.

\codeanchor{\raggedright The mesh is built by \fn{Fesom14Reader.load\_mesh} in \srcfile{readers/fesom14/reader.py},
with the helpers \fn{\_from\_ascii}, \fn{\_from\_polygons}, \fn{\_check\_deptho} and \fn{\_assemble}
of the same file; the ASCII files are parsed by \fn{read\_ascii\_mesh}
(\srcfile{readers/fesom14/\_ascii\_mesh.py}) and the triangles recovered by \fn{recover\_triangles}
(\srcfile{readers/fesom14/\_polygons.py});
the validity predicate \fn{is\_data} and \fn{wet\_mask\_levelwise} (\srcfile{readers/fesom14/\_layout.py});
the Euler rotation \fn{euler\_matrix} (\srcfile{core/rotation.py}).}


\section{Convex quadrilaterals and derived tables}\label{sec:mesh-derived}

The quadrilateral weights of \cref{sec:hvel-quad} have no fall-back, so every quadrilateral must be convex.
The FESOM-C reader, the structured-grid builder and the HEALPix dual builder check this when the mesh is
loaded (the HEALPix pixel builder too, except for the eight pixels that touch a pole, which it only reports
with a warning): each quadrilateral is projected onto the plane tangent to the sphere at its centre, and it fails
when the cross products of consecutive edges change sign or one is zero within a relative tolerance of
$10^{-9}$. The test does not depend on the winding direction or on a rotation of the frame, and a failure
names the offending elements.

Three further tables are built from the connectivity when a run needs them and are cached on the mesh:
\begin{description}[font=\normalfont\itshape,style=nextline,leftmargin=1.5em]
  \item[Element neighbours.] For each element, the element across each of its edges, or none at the coast;
    an edge shared by more than two elements is an error. Used by the edge least-squares scheme and by
    element tracking (\cref{sec:hvel-edge,sec:loc-walk}), and built when the run starts if either is on.
  \item[Coastal flags.] An element is coastal when one of its edges belongs to no other element; used by
    the beaching model (\cref{sec:life-beach}).
  \item[Element areas.] Read from the model's mesh file when present (FESOM2 mesh diagnostics provide them;
    ASCII meshes and regional extracts do not);
    otherwise the spherical excess of the element's fan of triangles, times the square of the mesh radius.
    Used only as ratios, in the element-to-vertex averages of the barycentric scheme.
\end{description}
The search structures of \cref{ch:localization} are built at the end of every loader.

\codeanchor{\fn{find\_nonconvex\_quads}, \fn{validate\_quad\_convexity}, \fn{build\_elem\_neighbors},
\fn{build\_boundary\_info} and \fn{build\_elem\_areas} in \srcfile{core/mesh.py}.}
